\PassOptionsToPackage{dvipsnames}{xcolor}
\documentclass[14pt,a4paper,oneside]{extreport}


\usepackage{cmap}
\usepackage[T2A]{fontenc}
\usepackage[utf8]{inputenc}
\usepackage[english,russian]{babel}
\usepackage{enumitem}
\usepackage{indentfirst}
\usepackage{amsmath,amsfonts,,amssymb,amsthm,mathtools}
\theoremstyle{definition}
\newtheorem{theorem}{Теорема}[section]
\newtheorem{corollary}[theorem]{Следствие}

\DeclareMathOperator{\DD}{d}

\usepackage{siunitx}
\sisetup{output-decimal-marker={,}}

\usepackage[version=4]{mhchem}

\usepackage{hyperref}
\hypersetup{
    colorlinks=true,
    linkcolor=blue!95!black,
    filecolor=black,
    urlcolor=black,
    citecolor=black,
    pdftitle={Курс современной школьной физики},
    pdfpagemode=UseOutlines,
}
\urlstyle{same}

\usepackage{graphicx,xcolor}
\usepackage{subcaption}
\graphicspath{figures}
\setlength\fboxsep{3pt}
\setlength\fboxrule{1pt}

\usepackage{float}
\usepackage{caption}
\captionsetup{
  font=small
}
\usepackage{anyfontsize}
\captionsetup{justification=centering}


\usepackage{geometry}
\geometry{
  a4paper,
  top=20mm,
  bottom=25mm,
  left=15mm,
  right=15mm
}
\newcommand{\name}[1]{%
  \begin{center}
    \vspace*{0.1cm}
    {\bfseries\fontsize{28}{32}\selectfont
      \parbox{\textwidth}{\centering #1\par}}%
  \end{center}%
}

\newcommand{\DEF}[2]{%
  \refstepcounter{definition}%
  \begin{DefinitionBox}
    \textbf{#1} — #2
  \end{DefinitionBox}
}

\newcommand{\opr}[3]{%
  \refstepcounter{definition}%
  \begin{DefinitionBox}
    \textbf{#1} $\equiv #2$ — #3
  \end{DefinitionBox}
}

\newcommand{\Opr}[4]{%
  \refstepcounter{definition}%
  \begin{DefinitionBox}
    \textbf{#1} $\equiv #2$ — #3; \quad $\left[#2\right] = #4$
  \end{DefinitionBox}
}

\newcommand{\Oprf}[4]{%
  \refstepcounter{definition}%
  \begin{DefinitionBox}
    \textbf{#1} $\equiv #2$ — #3;\\
    #4
  \end{DefinitionBox}
}


\newcommand{\DEFlabel}[3]{%
  \refstepcounter{definition}%
  \begin{DefinitionBox}[title={Определение \thedefinition}]%
    \textbf{#1} — #2%
    \label{#3}%
  \end{DefinitionBox}%
}

\newcommand{\zm}[1]{%
  \begin{RemarkBox}
    #1
  \end{RemarkBox}
}


\newcounter{definition}[chapter]
\renewcommand{\thedefinition}{\thechapter.\arabic{definition}}


\usepackage{tikz}
\usetikzlibrary{arrows.meta,calc,positioning,mindmap,shapes.symbols}

\usepackage{pgfplots}
\pgfplotsset{compat=newest}

\usepackage[most]{tcolorbox}
\tcbuselibrary{breakable, skins}

\newtcolorbox{DefinitionBox}[1][]{
    colback=brown!5!white,
    colframe=orange!80!black,
    fonttitle=\bfseries,
    title={Определение},
    breakable,
    enhanced,
    sharp corners=south,
    boxrule=0.5pt,
    #1
}


\newtcolorbox{RemarkBox}[1][]{
  colback=gray!10,
  colframe=gray,
  fonttitle=\bfseries,
  title={Замечание},
  #1,
  breakable
}

\newtcolorbox{ExampleBox}[1][]{
  colback=green!5!white,
  colframe=green!50!black,
  fonttitle=\bfseries,
  title={Пример},
  breakable,
  #1
}

\usetikzlibrary{decorations.markings,decorations.pathreplacing}

\begin{document}

\pagestyle{empty}
\pagenumbering{gobble}


\name{\S 6. Конденсатор}

Потенциальная энергия взаимодействия $W$ пробного заряда $q_0$ с электростатическим полем
в точке с потенциалом $\varphi$ определяется выражением
\begin{equation*}
  W = q_0\,\varphi.
\end{equation*}

В данном параграфе будет показано, что энергией обладает и само электростатическое поле. Для накопления и хранения этой энергии используются специальные устройства — \textit{конденсаторы.}
\subsection*{Электростатическое поле равномерно заряженной плоскости}

На рисунке \ref{fig:plate-1} показана картина линий напряжённости электростатического поля равномерно заряженной уединённой пластины.
Вблизи центральной части пластины линии напряжённости изображены \textit{сплошными}: в этой области поле с хорошей точностью можно считать \textit{однородным}.
Вблизи краёв пластины линии поля искривляются и изображены \textit{пунктиром}, что отражает \textit{неоднородность} электростатического поля в этих областях (\textit{краевые эффекты}).

На рисунке \ref{fig:plate-1a} представлен случай \textit{положительно} заряженной пластины, а на рисунке \ref{fig:plate-1b} — \textit{отрицательно} заряженной.
\begin{figure}[H]
  \centering

  \begin{subfigure}{0.48\textwidth}
    \centering
    \resizebox{\linewidth}{!}{%
      \begin{tikzpicture}[scale=1,>=latex]
      \def\ArrowDist{2cm}
      \tikzset{
        fieldarrow/.style={
          postaction={decorate},
          decoration={markings, mark=at position \ArrowDist with {\arrow{latex}}}
        }
      }

      \coordinate (A) at (-2.7,0);
      \coordinate (B) at (2.7,0);

      \foreach \x in {-2.5,-2,...,2.5} {
        \foreach \y in {-0.2,0.2}
      \node[red, font=\scriptsize] at (\x,\y) {$+$};
      }

      \foreach \x in {-0.25,0.25} {
        \foreach \y in {-3,3}
      \draw[semithick, MidnightBlue, postaction={decorate}, decoration={markings, mark=at position 0.85 with {\arrow{latex}}}] (\x,0) -- (\x,\y);
      }

      \foreach \x in {-0.75,0.75} {
        \foreach \y in {-3,3}
      \draw[semithick, MidnightBlue, postaction={decorate}, decoration={markings, mark=at position 0.85 with {\arrow{latex}}}] (\x,0) -- (\x,\y);
      }

      \draw[
        semithick,
        MidnightBlue
      ]
      (-1.25,0) .. controls (-1.25,0.9) and (-1.25,1.26) .. (-1.304,1.728);
      \draw[
        semithick,
        MidnightBlue,
        dashed,
        postaction={decorate},
        decoration={markings, mark=at position 0.5   with {\arrow{latex}}}
      ]
      (-1.304,1.728) .. controls (-1.34,2.04) and (-1.4,2.4) .. (-1.5,3);
      \draw[
        semithick,
        MidnightBlue
      ]
      (1.25,0) .. controls (1.25,0.9) and (1.25,1.26) .. (1.304,1.728);
      \draw[
        semithick,
        MidnightBlue,
        dashed,
        postaction={decorate},
        decoration={markings, mark=at position 0.5 with {\arrow{latex}}}
      ]
      (1.304,1.728) .. controls (1.34,2.04) and (1.4,2.4) .. (1.5,3);
      \draw[
        semithick,
        MidnightBlue
      ]
      (-1.25,0) .. controls (-1.25,-0.9) and (-1.25,-1.26) .. (-1.304,-1.728);
      \draw[
        semithick,
        MidnightBlue,
        dashed,
        postaction={decorate},
        decoration={markings, mark=at position 0.5 with {\arrow{latex}}}
      ]
      (-1.304,-1.728) .. controls (-1.34,-2.04) and (-1.4,-2.4) .. (-1.5,-3);
      \draw[
        semithick,
        MidnightBlue
      ]
      (1.25,0) .. controls (1.25,-0.9) and (1.25,-1.26) .. (1.304,-1.728);
      \draw[
        semithick,
        MidnightBlue,
        dashed,
        postaction={decorate},
        decoration={markings, mark=at position 0.5 with {\arrow{latex}}}
      ]
      (1.304,-1.728) .. controls (1.34,-2.04) and (1.4,-2.4) .. (1.5,-3);
      \draw[
        semithick,
        MidnightBlue
      ]
      (-1.75,0) .. controls (-1.79,0.6) and (-1.822,0.96) .. (-1.878,1.272);
      \draw[
        semithick,
        MidnightBlue,
        dashed,
        postaction={decorate},
        decoration={markings, mark=at position 0.5 with {\arrow{latex}}}
      ]
      (-1.878,1.272) .. controls (-1.962,1.74) and (-2.1,2.1) .. (-2.4,3);
      \draw[
        semithick,
        MidnightBlue
      ]
      (1.75,0) .. controls (1.79,0.6) and (1.822,0.96) .. (1.878,1.272);
      \draw[
        semithick,
        MidnightBlue,
        dashed,
        postaction={decorate},
        decoration={markings, mark=at position 0.5 with {\arrow{latex}}}
      ]
      (1.878,1.272) .. controls (1.962,1.74) and (2.1,2.1) .. (2.4,3);
      \draw[
        semithick,
        MidnightBlue
      ]
      (-1.75,0) .. controls (-1.79,-0.6) and (-1.822,-0.96) .. (-1.878,-1.272);
      \draw[
        semithick,
        MidnightBlue,
        dashed,
        postaction={decorate},
        decoration={markings, mark=at position 0.5 with {\arrow{latex}}}
      ]
      (-1.878,-1.272) .. controls (-1.962,-1.74) and (-2.1,-2.1) .. (-2.4,-3);
      \draw[
        semithick,
        MidnightBlue
      ]
      (1.75,0) .. controls (1.79,-0.6) and (1.822,-0.96) .. (1.8844,-1.272);
      \draw[
        semithick,
        MidnightBlue,
        dashed,
        postaction={decorate},
        decoration={markings, mark=at position 0.5 with {\arrow{latex}}}
      ]
      (1.8844,-1.272) .. controls (1.978,-1.74) and (2.14,-2.1) .. (2.5,-3);

      \draw[
        semithick,
        MidnightBlue
      ]
      (-2.25,0) .. controls (-2.265,0.45) and (-2.307,0.765) .. (-2.3868,1.026);
      \draw[
        semithick,
        MidnightBlue,
        dashed,
        postaction={decorate},
        decoration={markings, mark=at position 0.5 with {\arrow{latex}}}
      ]
      (-2.3868,1.026) .. controls (-2.573,1.635) and (-2.965,1.95) .. (-3.7,3);
      \draw[
        semithick,
        MidnightBlue
      ]
      (2.25,0) .. controls (2.265,0.45) and (2.307,0.765) .. (2.3868,1.026);
      \draw[
        semithick,
        MidnightBlue,
        dashed,
        postaction={decorate},
        decoration={markings, mark=at position 0.5 with {\arrow{latex}}}
      ]
      (2.3868,1.026) .. controls (2.573,1.635) and (2.965,1.95) .. (3.7,3);
      \draw[
        semithick,
        MidnightBlue
      ]
      (2.25,0) .. controls (2.265,-0.45) and (2.307,-0.765) .. (2.3868,-1.026);
      \draw[
        semithick,
        MidnightBlue,
        dashed,
        postaction={decorate},
        decoration={markings, mark=at position 0.5 with {\arrow{latex}}}
      ]
      (2.3868,-1.026) .. controls (2.573,-1.635) and (2.965,-1.95) .. (3.7,-3);
      \draw[
        semithick,
        MidnightBlue
      ]
      (-2.25,0) .. controls (-2.265,-0.45) and (-2.307,-0.765) .. (-2.3868,-1.026);
      \draw[
        semithick,
        MidnightBlue,
        dashed,
        postaction={decorate},
        decoration={markings, mark=at position 0.5 with {\arrow{latex}}}
      ]
      (-2.3868,-1.026) .. controls (-2.573,-1.635) and (-2.965,-1.95) .. (-3.7,-3);

      \draw[dashed,
        semithick,
        MidnightBlue,
        postaction={decorate},
        decoration={markings, mark=at position 0.65 with {\arrow{latex}}}
      ]
      (-2.5,0) -- (-3.7,0);
      \draw[dashed,
        semithick,
        MidnightBlue,
        postaction={decorate},
        decoration={markings, mark=at position 0.65 with {\arrow{latex}}}
      ]
      (2.5,0) -- (3.7,0);

      \draw[ultra thick] (A) -- (B);

      \node at (0,1.5) {$\vec{E}$};

      \end{tikzpicture}%
    }
    \caption{$q>0$.}
      \label{fig:plate-1a}
  \end{subfigure}\hspace{0.5cm}
  \begin{subfigure}{0.48\textwidth}
    \centering
    \resizebox{\linewidth}{!}{%
      \begin{tikzpicture}[scale=1,>=latex]
      \def\ArrowDist{2cm}
      \tikzset{
        fieldarrow/.style={
          postaction={decorate},
          decoration={markings, mark=at position \ArrowDist with {\arrow{Latex[reversed]}}}
        }
      }

      \coordinate (A) at (-2.7,0);
      \coordinate (B) at (2.7,0);

      \foreach \x in {-2.5,-2,...,2.5} {
        \foreach \y in {-0.2,0.2}
      \node[blue, font=\scriptsize] at (\x,\y) {$-$};
      }

      \foreach \x in {-0.25,0.25} {
        \foreach \y in {-3,3}
      \draw[semithick, MidnightBlue, postaction={decorate}, decoration={markings, mark=at position 0.85 with {\arrow{Latex[reversed]}}}] (\x,0) -- (\x,\y);
      }

      \foreach \x in {-0.75,0.75} {
        \foreach \y in {-3,3}
      \draw[semithick, MidnightBlue, postaction={decorate}, decoration={markings, mark=at position 0.85 with {\arrow{Latex[reversed]}}}] (\x,0) -- (\x,\y);
      }

      \draw[
        semithick,
        MidnightBlue
      ]
      (-1.25,0) .. controls (-1.25,0.9) and (-1.25,1.26) .. (-1.304,1.728);
      \draw[
        semithick,
        MidnightBlue,
        dashed,
        postaction={decorate},
        decoration={markings, mark=at position 0.5   with {\arrow{Latex[reversed]}}}
      ]
      (-1.304,1.728) .. controls (-1.34,2.04) and (-1.4,2.4) .. (-1.5,3);
      \draw[
        semithick,
        MidnightBlue
      ]
      (1.25,0) .. controls (1.25,0.9) and (1.25,1.26) .. (1.304,1.728);
      \draw[
        semithick,
        MidnightBlue,
        dashed,
        postaction={decorate},
        decoration={markings, mark=at position 0.5 with {\arrow{Latex[reversed]}}}
      ]
      (1.304,1.728) .. controls (1.34,2.04) and (1.4,2.4) .. (1.5,3);
      \draw[
        semithick,
        MidnightBlue
      ]
      (-1.25,0) .. controls (-1.25,-0.9) and (-1.25,-1.26) .. (-1.304,-1.728);
      \draw[
        semithick,
        MidnightBlue,
        dashed,
        postaction={decorate},
        decoration={markings, mark=at position 0.5 with {\arrow{Latex[reversed]}}}
      ]
      (-1.304,-1.728) .. controls (-1.34,-2.04) and (-1.4,-2.4) .. (-1.5,-3);
      \draw[
        semithick,
        MidnightBlue
      ]
      (1.25,0) .. controls (1.25,-0.9) and (1.25,-1.26) .. (1.304,-1.728);
      \draw[
        semithick,
        MidnightBlue,
        dashed,
        postaction={decorate},
        decoration={markings, mark=at position 0.5 with {\arrow{Latex[reversed]}}}
      ]
      (1.304,-1.728) .. controls (1.34,-2.04) and (1.4,-2.4) .. (1.5,-3);
      \draw[
        semithick,
        MidnightBlue
      ]
      (-1.75,0) .. controls (-1.79,0.6) and (-1.822,0.96) .. (-1.878,1.272);
      \draw[
        semithick,
        MidnightBlue,
        dashed,
        postaction={decorate},
        decoration={markings, mark=at position 0.5 with {\arrow{Latex[reversed]}}}
      ]
      (-1.878,1.272) .. controls (-1.962,1.74) and (-2.1,2.1) .. (-2.4,3);
      \draw[
        semithick,
        MidnightBlue
      ]
      (1.75,0) .. controls (1.79,0.6) and (1.822,0.96) .. (1.878,1.272);
      \draw[
        semithick,
        MidnightBlue,
        dashed,
        postaction={decorate},
        decoration={markings, mark=at position 0.5 with {\arrow{Latex[reversed]}}}
      ]
      (1.878,1.272) .. controls (1.962,1.74) and (2.1,2.1) .. (2.4,3);
      \draw[
        semithick,
        MidnightBlue
      ]
      (-1.75,0) .. controls (-1.79,-0.6) and (-1.822,-0.96) .. (-1.878,-1.272);
      \draw[
        semithick,
        MidnightBlue,
        dashed,
        postaction={decorate},
        decoration={markings, mark=at position 0.5 with {\arrow{Latex[reversed]}}}
      ]
      (-1.878,-1.272) .. controls (-1.962,-1.74) and (-2.1,-2.1) .. (-2.4,-3);
      \draw[
        semithick,
        MidnightBlue
      ]
      (1.75,0) .. controls (1.79,-0.6) and (1.822,-0.96) .. (1.8844,-1.272);
      \draw[
        semithick,
        MidnightBlue,
        dashed,
        postaction={decorate},
        decoration={markings, mark=at position 0.5 with {\arrow{Latex[reversed]}}}
      ]
      (1.8844,-1.272) .. controls (1.978,-1.74) and (2.14,-2.1) .. (2.5,-3);

      \draw[
        semithick,
        MidnightBlue
      ]
      (-2.25,0) .. controls (-2.265,0.45) and (-2.307,0.765) .. (-2.3868,1.026);
      \draw[
        semithick,
        MidnightBlue,
        dashed,
        postaction={decorate},
        decoration={markings, mark=at position 0.5 with {\arrow{Latex[reversed]}}}
      ]
      (-2.3868,1.026) .. controls (-2.573,1.635) and (-2.965,1.95) .. (-3.7,3);
      \draw[
        semithick,
        MidnightBlue
      ]
      (2.25,0) .. controls (2.265,0.45) and (2.307,0.765) .. (2.3868,1.026);
      \draw[
        semithick,
        MidnightBlue,
        dashed,
        postaction={decorate},
        decoration={markings, mark=at position 0.5 with {\arrow{Latex[reversed]}}}
      ]
      (2.3868,1.026) .. controls (2.573,1.635) and (2.965,1.95) .. (3.7,3);
      \draw[
        semithick,
        MidnightBlue
      ]
      (2.25,0) .. controls (2.265,-0.45) and (2.307,-0.765) .. (2.3868,-1.026);
      \draw[
        semithick,
        MidnightBlue,
        dashed,
        postaction={decorate},
        decoration={markings, mark=at position 0.5 with {\arrow{Latex[reversed]}}}
      ]
      (2.3868,-1.026) .. controls (2.573,-1.635) and (2.965,-1.95) .. (3.7,-3);
      \draw[
        semithick,
        MidnightBlue
      ]
      (-2.25,0) .. controls (-2.265,-0.45) and (-2.307,-0.765) .. (-2.3868,-1.026);
      \draw[
        semithick,
        MidnightBlue,
        dashed,
        postaction={decorate},
        decoration={markings, mark=at position 0.5 with {\arrow{Latex[reversed]}}}
      ]
      (-2.3868,-1.026) .. controls (-2.573,-1.635) and (-2.965,-1.95) .. (-3.7,-3);

      \draw[dashed,
        semithick,
        MidnightBlue,
        postaction={decorate},
        decoration={markings, mark=at position 0.65 with {\arrow{Latex[reversed]}}}
      ]
      (-2.5,0) -- (-3.7,0);
      \draw[dashed,
        semithick,
        MidnightBlue,
        postaction={decorate},
        decoration={markings, mark=at position 0.65 with {\arrow{Latex[reversed]}}}
      ]
      (2.5,0) -- (3.7,0);

      \draw[ultra thick] (A) -- (B);

      \node at (0,1.5) {$\vec{E}$};

      \end{tikzpicture}%
    }
    \caption{$q<0$.}
      \label{fig:plate-1b}
  \end{subfigure}
  \caption{Картина линий напряженности электростатиоческого поля  равномерно заряженной уединенной пластины. }
  \label{fig:plate-1}
\end{figure}

Используя \textit{теорему Гаусса}, можно показать, что в области, где электростатическое поле уединённой равномерно заряженной пластины можно считать \textit{однородным}, модуль напряжённости поля равен
\begin{equation}
E=\frac{|q|}{2\,\varepsilon\,\varepsilon_0\,S},
\label{eq:field-plate}
\end{equation}
где $q$ — заряд пластины, $S$ — её площадь, $\varepsilon$ — диэлектрическая проницаемость среды.

\subsection*{Плоский конденсатор}

Рассмотрим систему, изображённую на рисунке~\ref{fig:cond-1}, 
состоящую из двух параллельных пластин площадью $S$, 
имеющих противоположные заряды 
($+q$ — заряд положительной пластины, $-q$ — заряд отрицательной пластины), 
между которыми находится промежуток шириной $d$, 
заполненный диэлектриком с диэлектрической проницаемостью $\varepsilon$.
\begin{figure}[H]
  \centering
  \resizebox{0.28\linewidth}{!}{%
    \begin{tikzpicture}[scale=1,>=latex]
    \def\ArrowDist{2cm}
    \tikzset{
      fieldarrow/.style={
        postaction={decorate},
        decoration={markings, mark=at position \ArrowDist with {\arrow{latex}}}
      }
    }

    \coordinate (A) at (-1,-2.7);
    \coordinate (B) at (-1,2.7);
    \coordinate (C) at (1,-2.7);
    \coordinate (D) at (1,2.7);

    \fill[fill=gray, fill opacity=0.1] (A) rectangle (D);

    \foreach \x in {-1.2,-0.8} {
      \foreach \y in {-2.5,-2,...,2.5}
    \node[red, font=\scriptsize] at (\x,\y) {$+$};
    }
    \draw[ultra thick] (A) -- (B);

    \foreach \x in {1.2,0.8} {
      \foreach \y in {-2.5,-2,...,2.5}
    \node[blue, font=\scriptsize] at (\x,\y) {$-$};
    }
    \draw[ultra thick] (C) -- (D);

    \draw[CadetBlue,
      decorate,
      decoration={brace, raise=9pt},
    ] (-1,-2.7) -- (-1,2.7)
    node[midway, left=10pt] {$S$};
    \draw[CadetBlue,
      decorate,
      decoration={brace, mirror, raise=9pt},
    ] (1,-2.7) -- (1,2.7)
    node[midway, right=10pt] {$S$};

    \node[above] at (-1,2.7) {$+ q$};
    \node[above] at (1,2.7) {$- q$};

    \draw[CadetBlue,
      decorate,
      decoration={brace, mirror, raise=4pt},
    ] (-1,-2.7) -- (1,-2.7)
    node[midway, below=6pt] {$d$};

    \node at (0,0) {$\varepsilon$};

    \end{tikzpicture}%
  }
    \caption{Плоский конденсатор.}
      \label{fig:cond-1}

\end{figure}
Предполагается, что расстояние между пластинами значительно меньше размеров самих пластин ($d \ll \sqrt{S}$), вследствие чего краевыми эффектами можно пренебречь и электростатическое поле каждой пластины \textit{однородно}. Такую систему называют \textit{плоским конденсатором}. 

\DEF{Конденсатор}{система из двух проводников, несущих противоположные электрические заряды и разделённых тонким слоем диэлектрика.
}

\DEF{Плоский конденсатор}{
конденсатор, проводниками которого являются две параллельные плоские пластины, называемые обкладками.
}

\zm{Заряд конденсатора $q$ $\equiv$ заряд положительной обкладки $+q$ конденсатора.}

\subsubsection*{Напряженность поля конденсатора}

Каждая из обкладок конденсатора создаёт собственное электростатическое поле:
положительная обкладка — поле $\vec E_{+}$,
а отрицательная обкладка — поле $\vec E_{-}$,
как показано на рисунке~\ref{fig:cond-2}.

\begin{figure}[H]
  \centering
  \resizebox{0.49\linewidth}{!}{%
    \begin{tikzpicture}[scale=1,>=latex]
    \def\ArrowDist{2cm}
    \tikzset{
      fieldarrow/.style={
        postaction={decorate},
        decoration={markings, mark=at position \ArrowDist with {\arrow{latex}}}
      }
    }

    \coordinate (A) at (-1,-2.7);
    \coordinate (B) at (-1,2.7);
    \coordinate (C) at (1,-2.7);
    \coordinate (D) at (1,2.7);

    \fill[fill=gray, fill opacity=0.1] (A) rectangle (D);

    \foreach \x in {-1.2,-0.8} {
      \foreach \y in {-2.5,-2,...,2.5}
    \node[red, font=\scriptsize] at (\x,\y) {$+$};
    }
    \draw[ultra thick] (A) -- (B);

    \foreach \x in {1.2,0.8} {
      \foreach \y in {-2.5,-2,...,2.5}
    \node[blue, font=\scriptsize] at (\x,\y) {$-$};
    }
    \draw[ultra thick] (C) -- (D);

    \draw[CadetBlue,
      decorate,
      decoration={brace, raise=9pt},
    ] (-1,-2.7) -- (-1,2.7)
    node[midway, left=10pt] {$S$};
    \draw[CadetBlue,
      decorate,
      decoration={brace, mirror, raise=9pt},
    ] (1,-2.7) -- (1,2.7)
    node[midway, right=10pt] {$S$};

    \node[above] at (-1,2.7) {$+ q$};
    \node[above] at (1,2.7) {$- q$};

    \draw[CadetBlue,
      decorate,
      decoration={brace, mirror, raise=4pt},
    ] (-1,-2.7) -- (1,-2.7)
    node[midway, below=6pt] {$d$};

    \node at (0,0) {$\varepsilon$};

      \foreach \y in {-2.6,-1.6,...,2.4}
    \draw[->, red] (-2,\y) -- (-3,\y);
    \node[below, red] at (-2.5,2.45) {$\vec E_+$};

      \foreach \y in {-2.4,-1.4,...,2.6}
    \draw[<-, blue] (-2,\y) -- (-3,\y);
    \node[above=-3pt, blue] at (-2.5,2.65) {$\vec E_-$};

      \foreach \y in {-2.6,-1.6,...,2.4}
    \draw[->, red] (-0.5,\y) -- (0.5,\y);
    \node[below, red] at (0,2.45) {$\vec E_+$};

      \foreach \y in {-2.4,-1.4,...,2.6}
    \draw[->, blue] (-0.5,\y) -- (0.5,\y);
    \node[above=-3pt, blue] at (0,2.65) {$\vec E_-$};

      \foreach \y in {-2.6,-1.6,...,2.4}
    \draw[->, red] (2,\y) -- (3,\y);
    \node[below, red] at (2.5,2.45) {$\vec E_+$};

      \foreach \y in {-2.4,-1.4,...,2.6}
    \draw[<-, blue] (2,\y) -- (3,\y);
    \node[above=-3pt, blue] at (2.5,2.65) {$\vec E_-$};

    \end{tikzpicture}%
  }
    \caption{Напряженность поля каждой обкладки конденсатора: $\vec E_+$ --- поле $+q$, $\vec E_-$ --- поле $-q$.}
      \label{fig:cond-2}
\end{figure}

Модули напряженности $E_+$ и $E_-$ согласно выражению \ref{eq:field-plate} для электростатического поля каждой из равномерно заряженных пластин равны:
\begin{equation}
  E_+=E_-=\frac{q}{2\,\varepsilon\,\varepsilon_0\,S}. 
  \label{eq:fields-plates}
\end{equation}

Рассчитаем напряженность $\vec E$ поля конденсатора (результирующего электростатичского поля двух обкладок), используя принцип суперпозиции. 

В области пространства \textit{снаружи} конденсатора
\begin{equation}
   \vec E_\text{снаруж} = \vec E_+ + \vec E_-
\;\xrightarrow{\text{рис.}\ref{fig:cond-2}}\;
E_\text{снаруж} = |E_+ - E_-|
\;\xrightarrow{\eqref{eq:fields-plates}}\;
\begin{cases}
  E_\text{снаруж} = 0\\
  \vec E_\text{снаруж}=\vec 0.
\end{cases}
  \label{eq:field-plate-out}
\end{equation}

В области пространства \textit{внутри} конденсатора
\begin{equation}
   \vec E_\text{внутр} = \vec E_+ + \vec E_-
\;\xrightarrow{\text{рис.}\ref{fig:cond-2}}\;
E_\text{внутр} = E_+ + E_-
\;\xrightarrow{\eqref{eq:fields-plates}}\;
\begin{cases}
E_\text{внутр} = \frac{q}{\varepsilon\,\varepsilon_0\,S}\\
  \vec E_\text{внутр} \uparrow\uparrow \vec E_+.
\end{cases}
  \label{eq:field-plate-in}
\end{equation}

Таким образом, из результатов~\eqref{eq:field-plate-out} и~\eqref{eq:field-plate-in}
следует, что при пренебрежении краевыми эффектами \textit{однородное электростатическое поле
плоского конденсатора существует лишь в пространстве\allowbreak{} между его обкладками}. Полученный результат справедлив для конденсаторов произвольной формы
(в частности, сферических и цилиндрических).

\zm{Электростатическое поле конденсатора сосредоточено в пространстве между его обкладками. }

\subsubsection*{Емкость конденсатора}

На рисунке~\ref{fig:cond-3} изображена картина линий напряжённости
электростатического поля $\vec E$ \textit{идеального} плоского конденсатора,
то есть при условии $d \ll \sqrt{S}$.
Внутри \textit{реального} конденсатора силовые линии электростатического поля
имеют тот же общий вид, однако вблизи краёв обкладок возникают краевые эффекты: поле становится неоднородным, а линии напряжённости
искривляются. \textbf{Краевых эффекты учитывать на ЕГЭ не нужно.} 
\begin{figure}[H]
  \centering
  \resizebox{0.32\linewidth}{!}{%
    \begin{tikzpicture}[scale=1,>=latex]
    \def\ArrowDist{2cm}
    \tikzset{
      fieldarrow/.style={
        postaction={decorate},
        decoration={markings, mark=at position \ArrowDist with {\arrow{latex}}}
      }
    }

    \coordinate (A) at (-1,-2.7);
    \coordinate (B) at (-1,2.7);
    \coordinate (C) at (1,-2.7);
    \coordinate (D) at (1,2.7);

    \fill[fill=gray, fill opacity=0.1] (A) rectangle (D);

    \foreach \x in {-1.2,-0.8} {
      \foreach \y in {-2.5,-2,...,2.5}
    \node[red, font=\scriptsize] at (\x,\y) {$+$};
    }

      \foreach \y in {-2.25,-1.75,...,2.25}
    \draw[->, semithick, MidnightBlue] (-1,\y) -- (1,\y);
    \node[above, MidnightBlue] at (0,2.25) {$\vec E$};

    \draw[ultra thick] (A) -- (B);

    \foreach \x in {1.2,0.8} {
      \foreach \y in {-2.5,-2,...,2.5}
    \node[blue, font=\scriptsize] at (\x,\y) {$-$};
    }
    \draw[ultra thick] (C) -- (D);

    \draw[CadetBlue,
      decorate,
      decoration={brace, raise=9pt},
    ] (-1,-2.7) -- (-1,2.7)
    node[midway, left=10pt] {$S$};
    \draw[CadetBlue,
      decorate,
      decoration={brace, mirror, raise=9pt},
    ] (1,-2.7) -- (1,2.7)
    node[midway, right=10pt] {$S$};

    \node[above] at (-1,2.7) {$+ q$};
    \node[above] at (1,2.7) {$- q$};

    \draw[CadetBlue,
      decorate,
      decoration={brace, mirror, raise=4pt},
    ] (-1,-2.7) -- (1,-2.7)
    node[midway, below=6pt] {$d, U$};

    \node at (0,0) {$\varepsilon$};

    \end{tikzpicture}%
  }
    \caption{Картина линий напряженности $\vec E$ электростатического поля плоского конденсатора.}
      \label{fig:cond-3}
\end{figure}

\zm{Напряжение конденсатора $U \equiv$ напряжение между его обкладками. }

Поскольку электростатическое поле конденсатора однородно, напряжение его напряжение $U$ равно

\begin{equation}
U = E\,d \;\overset{\eqref{eq:field-plate-in}}{=}\;\frac{q\,d}{\varepsilon\,\varepsilon_0\,S}.
\label{eq:napr}
\end{equation}

Выразим из уравнения \ref{eq:napr} заряд плоского конденсатора $q$:
\begin{equation}
  q=\frac{\varepsilon\,\varepsilon_0\,S}{d}\,U.
  \label{eq:charge}
\end{equation} 

Согласно уравнению~\eqref{eq:charge}, $q \propto U$,
что справедливо для конденсаторов произвольной формы. При этом коэффициент
пропорциональности зависит от геометрии конденсатора и диэлектрических
свойств среды между его обкладками и называется
\textit{емкостью конденсатора}.
\opr{Емкость конденсатора}{C}{скалярная физическая величина, характеризующая способность конденсатора накапливать
электрический заряд и равная отношению заряда конденсатора к его
  напряжению:
\[
\begin{gathered}
  \boxed{C=\frac{q}{U}}\\
  \\
  [C]=\frac{\text{Кл}}{\text{В}}\equiv \text{Ф (фарад)}.
\end{gathered}
\]
} 

\zm{Емкость конденсатора $C$ численно равна заряду, сообщаемому конденсатору
для увеличения его напряжения на $1\ \text{В}$.}

\zm{Емкость конденсатора $C$ зависит только от его геометрии и диэлектрических свойств среды между его обкладками.}

Емкость $C$ плоского конденсатора равна
\begin{equation}
\boxed{
C=\frac{q}{U}
\;\overset{\eqref{eq:charge}}{=}\;
\frac{\varepsilon\,\varepsilon_0\,S}{d}
}
\label{eq:capacity}
\end{equation}

В частности, для емкости $C_0$ плоского \textit{воздушного} ($\varepsilon=1$) конденсатора
из уравнения~\eqref{eq:capacity} получаем выражение

\begin{equation}
\boxed{
C_0=\frac{\varepsilon_0\,S}{d}
}
\end{equation}

\subsubsection*{Энергия заряженного конденсатора}

Рассмотрим процесс зарядки исходно незаряженного плоского конденсатора с емкостью $C$ до заряда $q$, изображенный на рисунке \ref{fig:charging-capacitor}. 
\begin{figure}[H]
  \centering
  \begin{subfigure}[t]{0.28\textwidth}
    \centering
    \resizebox{\linewidth}{!}{%
      \begin{tikzpicture}[scale=1,>=latex]
      \useasboundingbox (-1.6,-3.4) rectangle (1.6,3.4);
      \def\ArrowDist{2cm}
      \tikzset{
        fieldarrow/.style={
          postaction={decorate},
          decoration={markings, mark=at position \ArrowDist with {\arrow{latex}}}
        }
      }

      \coordinate (A) at (-1,-2.7);
      \coordinate (B) at (-1,2.7);
      \coordinate (C) at (1,-2.7);
      \coordinate (D) at (1,2.7);

      \fill[fill=gray, fill opacity=0.1] (A) rectangle (D);

      \draw[ultra thick] (A) -- (B);

      \draw[ultra thick] (C) -- (D);

      \node[above] at (-1,2.7) {$0$};
      \node[above] at (1,2.7) {$0$};

      \end{tikzpicture}%
    }
    \caption{Незаряженный конденсатор}
    \label{fig:charging-capacitor-a}
  \end{subfigure}\hfill
  \begin{subfigure}[t]{0.28\textwidth}
    \centering
    \resizebox{\linewidth}{!}{%
      \begin{tikzpicture}[scale=1,>=latex]
      \useasboundingbox (-1.6,-3.4) rectangle (1.6,3.4);
      \def\ArrowDist{2cm}
      \tikzset{
        fieldarrow/.style={
          postaction={decorate},
          decoration={markings, mark=at position \ArrowDist with {\arrow{latex}}}
        }
      }

      \coordinate (A) at (-1,-2.7);
      \coordinate (B) at (-1,2.7);
      \coordinate (C) at (1,-2.7);
      \coordinate (D) at (1,2.7);

      \fill[fill=gray, fill opacity=0.1] (A) rectangle (D);

        \foreach \y in {-2,-1,1,2}
      \draw[->, MidnightBlue] (-1,\y) -- (1,\y);
      \node[above, MidnightBlue] at (0,2) {$\vec {\tilde E}$};

      \foreach \x in {-1.2,-0.8} {
        \foreach \y in {-2.5,-1.5,...,2.5}
      \node[red, font=\scriptsize] at (\x,\y) {$+$};
      }
      \draw[ultra thick] (A) -- (B);

      \foreach \x in {1.2,0.8} {
        \foreach \y in {-2.5,-1.5,...,2.5}
      \node[blue, font=\scriptsize] at (\x,\y) {$-$};
      }
      \draw[ultra thick] (C) -- (D);

      \draw[->, thick] (-0.5,0) -- (0.8,0)
        node[midway, above, font=\scriptsize] {$d\vec{F}_{\text{эл}}$};
      \filldraw[fill=red, draw=black] (-0.5,0) circle (0.08);
      \node[above, font=\scriptsize] at (-0.5,0) {$d\tilde q$};
      \draw[->,semithick] (0.9,-0.15) -- (-0.9,-0.15)
        node[midway,below,font=\scriptsize] {$\vec S$};

      \node[above] at (-1,2.7) {$+\tilde q$};
      \node[above] at (1,2.7) {$- \tilde q$};

      \draw[CadetBlue,
        decorate,
        decoration={brace, mirror, raise=4pt},
      ] (-1,-2.7) -- (1,-2.7)
      node[midway, below=6pt] {$\tilde U$};

      \end{tikzpicture}%
    }
    \caption{Перенос малого $\mathrm{d}\tilde q$ от пластины $-\tilde q$ к пластине $+\tilde q$}
    \label{fig:charging-capacitor-b}
  \end{subfigure}\hfill
  \begin{subfigure}[t]{0.28\textwidth}
    \centering
    \resizebox{\linewidth}{!}{%
      \begin{tikzpicture}[scale=1,>=latex]
      \useasboundingbox (-1.6,-3.4) rectangle (1.6,3.4);
      \def\ArrowDist{2cm}
      \tikzset{
        fieldarrow/.style={
          postaction={decorate},
          decoration={markings, mark=at position \ArrowDist with {\arrow{latex}}}
        }
      }

      \coordinate (A) at (-1,-2.7);
      \coordinate (B) at (-1,2.7);
      \coordinate (C) at (1,-2.7);
      \coordinate (D) at (1,2.7);

      \fill[fill=gray, fill opacity=0.1] (A) rectangle (D);

        \foreach \y in {-2.25,-1.75,...,2.25}
      \draw[->, MidnightBlue] (-1,\y) -- (1,\y);
      \node[above, MidnightBlue] at (0,2.25) {$\vec E$};

      \foreach \x in {-1.2,-0.8} {
        \foreach \y in {-2.5,-2,...,2.5}
      \node[red, font=\scriptsize] at (\x,\y) {$+$};
      }
      \draw[ultra thick] (A) -- (B);

      \foreach \x in {1.2,0.8} {
        \foreach \y in {-2.5,-2,...,2.5}
      \node[blue, font=\scriptsize] at (\x,\y) {$-$};
      }
      \draw[ultra thick] (C) -- (D);

      \node[above] at (-1,2.7) {$+ q$};
      \node[above] at (1,2.7) {$- q$};

      \draw[CadetBlue,
        decorate,
        decoration={brace, mirror, raise=4pt},
      ] (-1,-2.7) -- (1,-2.7)
      node[midway, below=6pt] {$U$};

      \end{tikzpicture}%
    }
    \caption{Заряженный конденсатор до заряда $q$}
    \label{fig:charging-capacitor-c}
  \end{subfigure}

  \caption{Процесс зарядки плоского конденсатора}
  \label{fig:charging-capacitor}
\end{figure}

\begin{enumerate}
  \item На рисунке~\ref{fig:charging-capacitor-a} изображён незаряженный конденсатор:
электростатическое поле отсутствует, напряжение конденсатора равно нулю.
\item На рисунке~\ref{fig:charging-capacitor-b} показан промежуточный момент
зарядки конденсатора: очередной малый положительный заряд $\mathrm{d}\tilde q$ переносится с отрицательной обкладки $-\tilde q$ на положительную обкладку $+\tilde q$. 
\item Наконец, на рисунке~\ref{fig:charging-capacitor-c} показано
конечное состояние: конденсатор заряжен до заряда $q$,
напряжение между его обкладками равно $U$. 
\end{enumerate}

Малая работа $\mathrm{d}A_{\text{эл}}$, совершаемая электрической силой
при переносе малого заряда (см. рис.~\ref{fig:charging-capacitor-b})
против электростатического поля, равна
\begin{equation}
\mathrm{d}A_{\text{эл}} = - \tilde U(\tilde q)\,\mathrm{d}\tilde q.
\label{eq:work}
\end{equation}
Знак «$-$» отражает тот факт, что при переносе заряда против
электростатического поля работа электрической силы отрицательна,
а $\tilde U(\tilde q)$ — текущее напряжение конденсатора
при его заряде $\tilde q$.

Текущее напряжение $\tilde U(\tilde q)$ выразим через заряд конденсатора,
используя определение его емкости:
\begin{equation}
  C = \frac{\tilde q}{\tilde U}
  \;\Rightarrow\;
  \tilde U(\tilde q) = \frac{\tilde q}{C}.
  \label{eq:napr-2}
\end{equation}

Подставляя выражение~\eqref{eq:napr-2} в~\eqref{eq:work}, получаем
\begin{equation}
  \mathrm{d}A_{\text{эл}} = - \frac{\tilde q}{C}\,\mathrm{d}\tilde q.
  \label{eq:work-2}
\end{equation}

Работа электрических сил $A_\text{эл}$ за весь процесс зарядки конденсатора до заряда $q$ складывается из малых работ $ \mathrm{d}A_{\text{эл}}$
\begin{equation}
A_{\text{эл}}
= \int\limits_{0}^{A_\text{эл}} \mathrm{d}A_{\text{эл}}
\overset{\eqref{eq:work-2}}{=}
- \int\limits_{0}^{q} \frac{\tilde q}{C}\,\mathrm{d}\tilde q=-\frac{q^2}{2C}
\label{eq:work-3}
\end{equation}

С другой стороны, электрическая сила является потенциальной.
Следовательно, её работа при зарядке конденсатора связана с изменением
потенциальной энергии системы взаимодействующих зарядов $W_C$:
\begin{equation}
  A_{\text{эл}} = W_{C,0} - W_C.
  \label{eq:work-energy}
\end{equation}

Здесь $W_{C,0}$ — потенциальная энергия системы при отсутствии зарядов
на обкладках конденсатора, а $W_C$ — потенциальная энергия взаимодействия
зарядов заряженного конденсатора.

Поскольку нулевой уровень потенциальной энергии выбирается произвольно,
положим
\begin{equation}
  W_{C,0} = 0.
  \label{eq:work-5}
\end{equation}

Потенциальная энергия $W_C$ обусловлена взаимодействием зарядов
на обкладках конденсатора, поэтому её принято называть
\textit{энергией заряженного конденсатора}.

Тогда из уравнений~\eqref{eq:work-3},~\eqref{eq:work-energy} и \eqref{eq:work-5} следует, что
энергия заряженного конденсатора равна
\begin{equation}
  W_C = \frac{q^2}{2C}.
  \label{eq:energy-1}
\end{equation}

Используя определение емкости конденсатора
\begin{equation}
C = \frac{q}{U},
\label{eq:cap-def}
\end{equation}
энергию заряженного конденсатора можно записать в различных эквивалентных формах.

Подставляя выражение~\eqref{eq:cap-def} в формулу~\eqref{eq:energy-1},
получаем
\begin{equation}
W_C = \frac{q\,U}{2}.
\label{eq:energy-2}
\end{equation}

Выражая заряд через напряжение, $q = C\,U$, из той же формулы~\eqref{eq:energy-1}
находим
\begin{equation}
W_C = \frac{C\,U^2}{2}.
\label{eq:energy-3}
\end{equation}

Таким образом, энергия заряженного конденсатора $W_C$ может быть выражена
любой из следующих эквивалентных формул:
\begin{equation}
\boxed{
W_C = \frac{q^2}{2C}
     = \frac{q\,U}{2}
     = \frac{C\,U^2}{2}.
}
\end{equation}

\subsection*{Энергия электрического поля}

Хотя энергия заряженного конденсатора $W_C$ формально выражается через заряд и напряжение,
физически она связана с электростатическим полем,
существующим между обкладками конденсатора.

В самом деле, электростатическое поле плоского конденсатора
локализовано в объёме между его обкладками.
Следовательно, энергия заряженного конденсатора
сосредоточена в этом объёме и может рассматриваться
как энергия электростатического поля.

Рассмотрим плоский конденсатор с площадью обкладок $S$
и расстоянием между ними $d$.
Объём области, в которой существует электростатическое поле, равен
\[
V = S\,d.
\]

Энергия заряженного конденсатора равна
\[
W_C = \frac{C\,U^2}{2}.
\]
Подставляя выражение для ёмкости плоского конденсатора
\[
C = \frac{\varepsilon\,\varepsilon_0\,S}{d},
\]
и учитывая, что напряжение между обкладками связано
с напряжённостью поля соотношением
\[
U = E\,d,
\]
получаем
\[
W_C
= \frac{\varepsilon\,\varepsilon_0\,S}{2d}\,(E\,d)^2
= \frac{\varepsilon\,\varepsilon_0}{2}\,E^2\,S\,d.
\]

Таким образом, энергия электростатического поля,
содержащегося в объёме $V = S\,d$, равна
\[
W_C = \frac{\varepsilon\,\varepsilon_0}{2}\,E^2\,V.
\]

Отсюда следует, что \textit{плотность энергии электростатического поля}
равна
\begin{equation}
  \boxed{
  w = \frac{W_C}{V}
  = \frac{\varepsilon\,\varepsilon_0}{2}\,E^2 }
\end{equation}

Полученное выражение показывает,
что энергия электростатического поля является локальной величиной
и в каждой точке пространства определяется значением напряжённости поля.

Важно, что данный результат остаётся справедливым
для произвольного электростатического поля,
а также, в более общем случае, для электрических полей в целом.


\end{document}
